\documentclass[11pt,a4paper]{article}

\usepackage[margin=2.4cm]{geometry}
\usepackage[T1]{fontenc}
\usepackage[utf8]{inputenc}
\usepackage{lmodern}
\usepackage{amsmath,amssymb}
\usepackage{booktabs}
\usepackage{microtype}
\usepackage{enumitem}
\usepackage[svgnames]{xcolor}
\usepackage{upquote}
\usepackage[hidelinks]{hyperref}

\hypersetup{
  colorlinks=true,
  linkcolor=NavyBlue,
  urlcolor=NavyBlue,
  citecolor=NavyBlue,
  pdftitle={BTPC: Bitcoin-Time Protocol Chain — White Paper},
  pdfauthor={BTPC Development Team},
  pdfsubject={A quantum-resistant proof-of-work cryptocurrency with sustainable economics}
}

\title{\textbf{BTPC: Bitcoin-Time Protocol Chain}\\[0.6em]
\large A Quantum-Resistant Proof-of-Work Cryptocurrency with Sustainable Economics}
\author{BTPC Development Team}
\date{Version 2.2 --- November 2026}

\setlength{\parskip}{0.5em}
\setlength{\parindent}{0pt}

\begin{document}

\maketitle

\begin{abstract}
\noindent BTPC (Bitcoin-Time Protocol Chain) is a quantum-resistant, proof-of-work cryptocurrency that addresses two structural weaknesses of Bitcoin-class systems: vulnerability to cryptographically relevant quantum computers, and the long-term economic fragility of geometric halving schedules.

BTPC signs transactions with ML-DSA (Dilithium5, NIST post-quantum signature level~5), hashes with SHA-512 (64-byte digests), and derives keys with SHAKE-256 under domain separation. Its monetary policy replaces Bitcoin's stepwise halving with a continuous linear decay: the block reward starts at 32.375~BTPC and decreases linearly (an average of $\approx$2{,}525 credits per block) over 24 years (1{,}262{,}304 blocks), reaching a permanent tail emission of 0.5~BTPC. Total supply at the end of the decay period is approximately 20.75~million BTPC, after which annual issuance is fixed at $\sim$26{,}298~BTPC and the inflation rate asymptotically approaches zero.

The core is a clean-room Rust implementation following Bitcoin's proven architecture --- the UTXO model, proof-of-work consensus, 10-minute target block times, and 2{,}016-block difficulty adjustment --- with every parameter verified against the codebase (see Appendix~A).
\end{abstract}

\tableofcontents
\newpage

\section{Introduction}

\subsection{The Quantum Threat}

Bitcoin and nearly all deployed blockchains secure ownership with ECDSA over secp256k1. Shor's algorithm breaks the elliptic-curve discrete-logarithm problem in polynomial time on a sufficiently large fault-tolerant quantum computer, allowing an adversary to recover private keys from exposed public keys and forge signatures. Grover's algorithm halves the effective security of symmetric primitives: SHA-256's 256-bit classical security drops to 128-bit against quantum search.

Blockchain transactions are permanent and publicly visible. The ``harvest now, decrypt later'' model means funds secured by quantum-vulnerable signatures are exposed today for attack the moment a capable machine arrives. BTPC's design premise is that long-term value storage cannot wait for that event.

\subsection{The Halving Problem}

Bitcoin's halving reduces the block subsidy by 50\% every 210{,}000 blocks ($\sim$4 years). This creates abrupt, discontinuous revenue shocks for miners, transient hash-rate instability after each event, and --- over the long run --- a transition to a fee-only security model whose viability is unproven. No major blockchain has yet operated at scale under a fee-dominant security regime.

\subsection{BTPC's Approach}

BTPC addresses both problems in a single coherent design: post-quantum cryptography throughout, and a smooth linear-decay reward schedule that removes halving shocks while guaranteeing a perpetual minimum miner subsidy.

\section{Design Principles}

\begin{itemize}[leftmargin=1.5em,itemsep=2pt]
\item \textbf{Quantum-first.} Every cryptographic primitive is post-quantum. No classical public-key cryptography (RSA, ECDSA, DSA) appears anywhere in the consensus-critical path.
\item \textbf{Bitcoin-compatible architecture.} BTPC deliberately retains the structures that have survived 15 years of adversarial testing: the UTXO model, SHA-based proof-of-work, 10-minute target block times, and 2{,}016-block difficulty adjustment.
\item \textbf{Predictable economics.} The emission schedule is a closed-form arithmetic function of block height. Any participant can compute the exact reward at any future height with integer arithmetic --- no governance votes, oracles, or parameter updates.
\item \textbf{Security through simplicity.} BTPC supports only Pay-to-Public-Key-Hash (P2PKH) transactions. There is no scripting language and no Turing-complete execution environment.
\item \textbf{Memory safety.} The core is written in Rust, eliminating buffer overflows, use-after-free, and data races at the language level.
\end{itemize}

\section{Cryptographic Foundations}

\subsection{Signatures: ML-DSA (Dilithium5)}

BTPC uses ML-DSA at NIST security level~5 (the Dilithium5 parameter set). It is the only signature scheme used for transaction authorisation.

\begin{center}
\begin{tabular}{lr}
\toprule
\textbf{Parameter} & \textbf{Size} \\
\midrule
Private key & 4{,}864 bytes \\
Public key & 2{,}592 bytes \\
Signature & 4{,}595 bytes \\
\bottomrule
\end{tabular}
\end{center}

The security basis is the Module-Learning-With-Errors (MLWE) and Module-Short-Integer-Solution (MSIS) lattice problems. The large signatures and keys are the unavoidable cost of post-quantum security; BTPC accepts this in exchange for eliminating the entire class of quantum signature-forgery attacks.

\subsection{Hashing: SHA-512}

All hashing uses SHA-512, producing 64-byte digests. Block hashes are double-hashed, following Bitcoin's pattern:

\begin{verbatim}
block_hash = SHA-512(SHA-512(block_header))
\end{verbatim}

SHA-512 provides 256-bit security against Grover's algorithm (versus 128-bit for SHA-256), and the double hash prevents length-extension attacks.

\subsection{Key Derivation: SHAKE-256}

Wallet seed expansion uses SHAKE-256 (the SHA-3 extendable-output function) with a domain-separation tag so that the same BIP-39 mnemonic in another system yields different keys. The pipeline is: 32-byte entropy $\to$ BIP-39 mnemonic $\to$ SHAKE-256 expansion (domain-separated) $\to$ ML-DSA key generation. All-zero seeds are rejected.

\subsection{Wallet Encryption}

Wallet files are encrypted at rest with Argon2id (memory-hard password hashing, 64~MB / 3 iterations / 4 lanes) feeding AES-256-GCM authenticated encryption. Sensitive material (passwords, seeds, private keys) is held in zeroizing containers that wipe memory on drop.

\subsection{Constant-Time Operations}

Difficulty-target comparisons and credential verification use constant-time equality checks to avoid timing side channels.

\section{Address Format}

BTPC addresses use Base58Check encoding with a network version byte, mirroring Bitcoin:

\begin{center}
\begin{tabular}{lcl}
\toprule
\textbf{Network} & \textbf{Version byte} & \textbf{Prefix} \\
\midrule
Mainnet & \texttt{0x00} & \texttt{1} \\
Testnet & \texttt{0x6f} & \texttt{m}/\texttt{n} \\
Regtest & \texttt{0x6f} & \texttt{m}/\texttt{n} \\
\bottomrule
\end{tabular}
\end{center}

An address is derived from the public key as:

\begin{verbatim}
address = Base58Check(version_byte || RIPEMD-160(SHA-512(public_key)))
\end{verbatim}

Note the double SHA-512 checksum (rather than Bitcoin's double SHA-256) and SHA-512 (rather than SHA-256) in the address hash, matching the 64-byte hash space used throughout.

\section{Economic Model: Linear Decay}

\subsection{Constants}

\begin{center}
\begin{tabular}{ll}
\toprule
\textbf{Constant} & \textbf{Value} \\
\midrule
Initial reward ($R_0$) & 3{,}237{,}500{,}000 credits (32.375 BTPC) \\
Tail emission ($R_T$) & 50{,}000{,}000 credits (0.5 BTPC) \\
Blocks per year ($B$) & 52{,}596 ($365.25 \times 24 \times 6$) \\
Decay period ($Y$) & 24 years \\
Decay end height ($N$) & 1{,}262{,}304 blocks \\
Average decrease per block & $\approx$ 2{,}525.1445 credits (floor 2{,}525) \\
Base unit & 1 BTPC $=$ 100{,}000{,}000 credits \\
\bottomrule
\end{tabular}
\end{center}

\subsection{Reward Function}

Let $D = R_0 - R_T = 3{,}187{,}500{,}000$ be the total decay amount and $N = Y \cdot B = 1{,}262{,}304$ the decay-end height.

\[
R(h) =
\begin{cases}
R_0, & h = 0 \\
\max\bigl(R_0 - \lfloor D \cdot h / N \rfloor,\; R_T\bigr), & 0 < h < N \\
R_T, & h \ge N
\end{cases}
\]

The reward is computed with integer-only arithmetic (u128 intermediates) so the result is bit-identical on every platform --- a consensus requirement. The average decrease rate is $\Delta = D/N \approx 2525.1445$ credits per block; the stored constant \texttt{DECREASE\_PER\_BLOCK} $= \lfloor \Delta \rfloor = 2525$ is its floor.

\subsection{Reward Schedule}

\begin{center}
\begin{tabular}{lrr}
\toprule
\textbf{Height} & \textbf{Reward (credits)} & \textbf{Reward (BTPC)} \\
\midrule
0 (genesis) & 3{,}237{,}500{,}000 & 32.375 \\
52{,}596 (year 1) & 3{,}104{,}687{,}500 & 31.047 \\
631{,}152 (year 12) & 1{,}643{,}750{,}000 & 16.438 \\
1{,}262{,}303 (last decay block) & 50{,}002{,}526 & 0.500 \\
1{,}262{,}304 (first tail block) & 50{,}000{,}000 & 0.500 \\
\bottomrule
\end{tabular}
\end{center}

\subsection{Supply and Inflation}

Total supply at the end of the decay period (year 24) is approximately \textbf{20.75~million BTPC}. Thereafter, issuance is a fixed $\sim$26{,}298~BTPC per year ($0.5 \times 52{,}596$ blocks), and the inflation \emph{rate} declines monotonically toward zero as the supply grows. The perpetual tail emission provides a permanent miner subsidy independent of fee-market conditions, and gradually replaces coins lost to key loss and user error.

\section{Consensus Mechanism}

\subsection{Proof of Work}

Proof-of-work is SHA-512 based. A block is valid when its double-SHA-512 header hash is numerically below the current difficulty target. The nonce is 4 bytes.

\subsection{Difficulty Adjustment}

Difficulty is recalculated every 2{,}016 blocks with a target timespan of 1{,}209{,}600 seconds (two weeks), retargeting toward the 600-second block interval. This preserves Bitcoin's proven responsiveness while operating in the 64-byte hash space.

\subsection{Timestamp Rules}

\begin{itemize}[leftmargin=1.5em,itemsep=2pt]
\item Target block time: 600 seconds.
\item Minimum spacing between consecutive blocks: 60 seconds.
\item Maximum future timestamp: 7{,}200 seconds (2 hours) ahead of network time.
\item Timestamps are validated against median-time-past, Bitcoin-style.
\end{itemize}

\section{Transaction System}

BTPC uses the UTXO model with Bitcoin-compatible structure. Each transaction consumes one or more outputs and creates new ones; ownership is transferred by satisfying a P2PKH script with a valid ML-DSA signature.

Only P2PKH is supported. There are no smart contracts, no multi-signature aggregation, and no scripting beyond the minimal \texttt{OP\_DUP OP\_HASH160 \textless{}hash\textgreater{} OP\_EQUALVERIFY OP\_CHECKSIG} pattern. A transaction carries a \texttt{fork\_id} byte identifying the network (mainnet / testnet / regtest) to prevent cross-network replay.

\section{Block Structure}

The block header mirrors Bitcoin's, with the two hash fields expanded to 64 bytes for SHA-512:

\begin{center}
\begin{tabular}{lr}
\toprule
\textbf{Field} & \textbf{Size} \\
\midrule
Version & 4 bytes \\
Previous block hash & 64 bytes \\
Merkle root & 64 bytes \\
Timestamp & 4 bytes \\
Bits (difficulty) & 4 bytes \\
Nonce & 4 bytes \\
\midrule
\textbf{Total} & \textbf{144 bytes} \\
\bottomrule
\end{tabular}
\end{center}

The 144-byte header (versus Bitcoin's 80 bytes) is entirely attributable to the 64-byte digests replacing 32-byte ones.

\section{Network Architecture}

\subsection{Network Parameters}

\begin{center}
\begin{tabular}{lccc}
\toprule
\textbf{Network} & \textbf{P2P port} & \textbf{RPC port} & \textbf{Magic bytes} \\
\midrule
Mainnet & 18341 & 18340 & \texttt{B7 C0 1A 55} \\
Testnet & 18351 & 18350 & \texttt{F1 C0 BA 55} \\
Regtest & 18361 & 18360 & \texttt{DA B5 BF FA} \\
\bottomrule
\end{tabular}
\end{center}

Each network has a distinct P2P port and 4-byte magic value, preventing accidental cross-network connections. Nodes listen dual-stack on \texttt{[::]:<port>}.

\subsection{Peer Discovery and Synchronisation}

The network follows Bitcoin's bootstrap model:

\begin{itemize}[leftmargin=1.5em,itemsep=2pt]
\item A persistent address manager (the analogue of Bitcoin's \texttt{addrman}) stores known peers in new/tried tables, bucketed by /16 (IPv4) and /32 (IPv6) network groups, and survives restarts via RocksDB.
\item Bootstrap proceeds from a persisted \texttt{peers.dat} address book, LAN UDP discovery, DNS seeds, and manual peer-add --- in that order. (At present the network is behind carrier-grade NAT with no hardcoded seed nodes, so fresh nodes rely on the address book and manual peer-add until a reachable seed is established.)
\item Headers-first synchronisation is used, with per-peer and aggregate in-flight request caps, and a stall reaper that disconnects and penalises unresponsive peers.
\item Self-connections are detected by comparing a random nonce exchanged in the version handshake, preventing a node from counting itself as a peer.
\end{itemize}

\subsection{RPC Interface}

Nodes expose an authenticated JSON-RPC 2.0 interface (mainnet 18340, testnet 18350, regtest 18360) used by the mining rig and operator tooling. Authentication is mandatory --- the node refuses to start without valid credentials --- and the interface can be secured with TLS (a self-signed certificate is generated by the one-command \texttt{--setup} flow).

\section{Security Considerations}

\begin{itemize}[leftmargin=1.5em,itemsep=2pt]
\item \textbf{Post-quantum signatures} eliminate quantum forgery and key-recovery attacks against transaction ownership.
\item \textbf{Constant-time operations} mitigate timing side channels in cryptographic comparisons.
\item \textbf{Memory-safe implementation} (Rust) removes entire vulnerability classes.
\item \textbf{Minimal scripting surface} (P2PKH only) sharply reduces the attack surface relative to general-purpose platforms.
\item \textbf{Rate limiting and connection caps} (per-IP and per-subnet) mitigate denial-of-service and eclipse attacks.
\item \textbf{Authenticated, TLS-capable RPC} protects the operator interface.
\end{itemize}

Known trade-offs, stated honestly: ML-DSA signatures and keys are orders of magnitude larger than ECDSA (4{,}595-byte signatures vs $\sim$71 bytes), which increases on-chain storage and bandwidth costs; there is no native multi-signature aggregation; and ML-DSA is not yet supported in commercial hardware wallets.

\section{Comparison with Bitcoin}

\begin{center}
\begin{tabular}{lll}
\toprule
\textbf{Property} & \textbf{Bitcoin} & \textbf{BTPC} \\
\midrule
Signature scheme & ECDSA (secp256k1) & ML-DSA (Dilithium5) \\
Quantum resistance & No & Yes (NIST level 5) \\
Hash function & SHA-256 (32 bytes) & SHA-512 (64 bytes) \\
Header size & 80 bytes & 144 bytes \\
Block time & $\sim$10 minutes & $\sim$10 minutes \\
Difficulty period & 2{,}016 blocks & 2{,}016 blocks \\
Reward model & Geometric halving & Linear decay + tail emission \\
Scripting & Full script & P2PKH only \\
Base unit & satoshi & credit \\
\bottomrule
\end{tabular}
\end{center}

\section{Implementation Status}

The core is implemented in Rust as a workspace: \texttt{btpc-core} (the library), \texttt{btpc\_node} (headless full node), \texttt{btpc\_wallet} (CLI wallet), \texttt{btpc\_miner} (CPU miner), and a Tauri desktop wallet. A dedicated GPU mining rig (\texttt{BTPC\_RIG}) mines against a node over authenticated HTTPS, and a BalenaEtcher-flashable OS image (\texttt{BTPC-OS}) boots commodity hardware straight into mining with NVIDIA and AMD GPU support.

Post-quantum cryptography, the linear-decay consensus, P2P networking with Bitcoin-style peer discovery, and the authenticated RPC layer are implemented and covered by automated tests. The project is near completion, with peer-discovery hardening and node--rig security (mandatory auth + TLS + certificate pinning) being the most recent work.

\section{Conclusion}

BTPC demonstrates that quantum resistance and economic sustainability are achievable without abandoning the architectural lessons of Bitcoin. By pairing NIST-level post-quantum signatures and SHA-512 hashing with a continuous linear-decay monetary policy and a permanent tail emission, BTPC provides a store of value designed for the multi-decade horizon in which quantum computers will exist.

\appendix

\section{Parameter Reference}

Every protocol parameter in this document is verified against the implementation. Source locations (live tree, \texttt{/media/bob/HDD1/BTPC}):

\begin{center}
\begin{tabular}{llp{6.4cm}}
\toprule
\textbf{Parameter} & \textbf{Value} & \textbf{Source} \\
\midrule
P2P ports (main/test/reg) & 18341 / 18351 / 18361 & \texttt{btpc-core/src/network/mod.rs} \\
Magic bytes & \texttt{B7C01A55}/\texttt{F1C0BA55}/\texttt{DAB5BFFA} & \texttt{btpc-core/src/network/protocol.rs} \\
RPC ports (main/test/reg) & 18340 / 18350 / 18360 & \texttt{btpc-core/src/economics/constants.rs} \\
Protocol version & 70015 & \texttt{btpc-core/src/network/protocol.rs} \\
Target block time & 600 s & \texttt{btpc-core/src/consensus/mod.rs} \\
Min block spacing & 60 s & \texttt{btpc-core/src/consensus/mod.rs} \\
Max future block time & 7200 s & \texttt{btpc-core/src/consensus/mod.rs} \\
Difficulty period & 2{,}016 blocks & \texttt{btpc-core/src/consensus/difficulty.rs} \\
Genesis timestamp & 1735689600 (2025-01-01 UTC) & \texttt{btpc-core/src/blockchain/genesis.rs} \\
Initial reward & 3{,}237{,}500{,}000 credits & \texttt{btpc-core/src/consensus/rewards.rs} \\
Tail emission & 50{,}000{,}000 credits & (ibid.) \\
Decay end height & 1{,}262{,}304 blocks & (ibid.) \\
Average decrease per block & $\approx$ 2{,}525.1445 credits & \texttt{btpc-core/src/economics/constants.rs} \\
ML-DSA private key & 4{,}864 bytes & \texttt{btpc-core/src/crypto/mod.rs} \\
ML-DSA public key & 2{,}592 bytes & \texttt{btpc-core/src/crypto/mod.rs} \\
ML-DSA signature & 4{,}595 bytes & \texttt{btpc-core/src/crypto/mod.rs} \\
Address version bytes & 0x00 / 0x6f & \texttt{btpc-core/src/crypto/address.rs} \\
Credits per BTPC & 100{,}000{,}000 & economics module \\
\bottomrule
\end{tabular}
\end{center}

The economic model is independently verified by the \texttt{BTPC\_VERIFICATION\_ALGORITHM} script, which asserts the reward function at key heights, the decay-end supply (between 20.5~M and 21~M BTPC), and the year-over-year inflation schedule.

\section{Mathematical Specification}

This appendix states the exact consensus and cryptographic calculations, matching the implementation. Let $H(x) = \text{SHA-512}(x)$, and let $\lfloor \cdot \rfloor$ denote integer (floor) division.

\subsection{Hashing}

Block and transaction identifiers use double SHA-512:

\[
\operatorname{id}(x) = H(H(x))
\]

\subsection{Block Reward (Linear Decay)}

With $R_0 = 3{,}237{,}500{,}000$, $R_T = 50{,}000{,}000$, $B = 52{,}596$, $Y = 24$, $N = YB = 1{,}262{,}304$, and $D = R_0 - R_T = 3{,}187{,}500{,}000$:

\[
R(h) =
\begin{cases}
R_0, & h = 0 \\
\max\bigl(R_0 - \lfloor D \cdot h / N \rfloor,\; R_T\bigr), & 0 < h < N \\
R_T, & h \ge N
\end{cases}
\]

The per-block decrease is $\Delta = D/N \approx 2525.1445$ credits/block; the stored constant \texttt{DECREASE\_PER\_BLOCK} $= \lfloor \Delta \rfloor = 2525$ is its floor.

\subsection{Total Supply}

\[
S(h) = \sum_{i=0}^{h} R(i)
\]

Decay-period closed form (arithmetic series, used by \texttt{MAX\_MONEY}):

\[
S(N) \approx \frac{N (R_0 + R_T)}{2} = \frac{1{,}262{,}304 \times 3{,}287{,}500{,}000}{2} \approx 2.075 \times 10^{15}\ \text{credits} \approx 20.75\ \text{M BTPC}
\]

\subsection{Inflation Rate}

\[
I(h) = \frac{R(h) \cdot B}{S(h)} \times 100\%
\]

\subsection{Maximum Supply}

\[
\mathit{MAX\_MONEY} = \frac{N (R_0 + R_T)}{2} + 1000 \cdot B \cdot R_T
\]

(the tail term uses 1000 years as the practical horizon).

\subsection{Difficulty Retarget}

Every $N_d = 2016$ blocks, with target timespan $T_{\text{target}} = 2016 \times 600\ \text{s} = 1{,}209{,}600\ \text{s}$:

\[
r = \frac{T_{\text{actual}}}{T_{\text{target}}}, \qquad
r = \operatorname{clamp}\!\left(r,\ \tfrac{1}{4},\ 4\right), \qquad
T_{\text{next}} = T_{\text{prev}} \times r
\]

\subsection{Proof of Work}

A block is valid if and only if its double-SHA-512 header hash is at or below the target (64-byte big-endian comparison):

\[
H(H(\text{header})) \le \text{target}
\]

\subsection{Compact Target Encoding (bits)}

The 64-byte target is serialised as a 32-bit compact ``bits'' value --- a 1-byte exponent plus a 3-byte mantissa, Bitcoin-style, with the mantissa occupying the most-significant bytes of the 64-byte field.

\subsection{Merkle Root}

Leaves are transaction IDs. Each internal node is the double SHA-512 of the concatenation of its two children; an odd node is duplicated (Bitcoin behaviour):

\[
\text{node} = H\bigl(H(\text{left} \mathbin{\Vert} \text{right})\bigr)
\]

\subsection{Transaction Fee}

\[
\text{fee} = \sum \text{input\_UTXO\_values} - \sum \text{output\_values}
\]

\subsection{Address Derivation}

\[
\text{address} = \text{Base58Check}\bigl(\text{version} \mathbin{\Vert} \text{RIPEMD-160}(H(\text{public\_key}))\bigr)
\]
\[
\text{checksum} = H(H(\text{payload}))[0..4]
\]

The Base58Check checksum is the first 4 bytes of the double SHA-512 of the payload.

\end{document}
